\section{信息率失真函数}

\begin{remarkBox}{本章重点}
本章研究有失真信源编码的理论下界。学习主线是：先用失真测度描述原始符号 $X$ 与再现符号 $Y$ 之间的误差，再用平均失真 $D$ 约束允许的复现质量；随后定义信息率失真函数 $R(D)$，理解它是在给定失真约束下所需的最小互信息或最小编码率；最后掌握限失真信源编码定理和限失真信源信道编码定理，知道何时能以不超过 $D$ 的平均失真完成传输。
\end{remarkBox}

前面第五章讨论无失真信源编码，它要求译码后完全恢复信源序列，因此码率下界由熵或熵率给出。本章放宽要求：只要复现序列与原序列的平均差异不超过给定限度，就允许压缩得更低。这样得到的理论边界就是信息率失真函数，也称率失真函数。

理解本章时要分清两个方向：信道容量是在信道给定时最大化 $I(X;Y)$，回答“信道最多能传多少”；率失真理论是在信源和失真准则给定时最小化 $I(X;Y)$，回答“达到某个复现质量至少要传多少”。

\subsection{限失真信源编码模型}

有失真压缩系统由信源、信源编码器、无噪信道、信源译码器和信宿组成。信源输出消息 $X$，编码器把它压缩成较短的码字或码字序号，经无噪信道送到译码器，译码器输出复现消息 $Y$。由于压缩允许失真，通常 $Y$ 不再是 $X$ 的精确复制。

从信息论角度看，从 $X$ 到 $Y$ 的整体作用可以抽象为一个条件概率分布 $p(y\mid x)$。这个分布不一定是真实物理信道，而是描述“给定信源符号 $x$ 后，复现符号 $y$ 以多大概率出现”的\textbf{试验信道}。率失真函数就是在所有满足失真约束的试验信道中，寻找互信息最小的一个。

视频直播中的“超清但卡顿”和“标清但流畅”可以作为直观例子：如果信道或存储资源有限，完全无失真地传输高质量视频可能代价过高；适当降低画质可以显著降低码率，只要失真仍在用户可接受范围内。

\subsection{失真测度与平均失真}

\subsubsection{单符号失真测度}

设信源符号集为
\[
A=\{a_1,a_2,\dots,a_n\},
\]
再现符号集为
\[
B=\{b_1,b_2,\dots,b_m\}.
\]
当输入为 $x\in A$、输出为 $y\in B$ 时，用
\[
d(x,y)\ge0
\]
表示二者之间的失真或代价。失真越小，说明复现越接近原符号。

\begin{definitionBox}{失真矩阵}
单符号失真测度可以写成失真矩阵
\[
\mathbf d=
\begin{bmatrix}
d(a_1,b_1)&d(a_1,b_2)&\cdots&d(a_1,b_m)\\
d(a_2,b_1)&d(a_2,b_2)&\cdots&d(a_2,b_m)\\
\vdots&\vdots&\ddots&\vdots\\
d(a_n,b_1)&d(a_n,b_2)&\cdots&d(a_n,b_m)
\end{bmatrix}.
\]
其中第 $i$ 行第 $j$ 列元素 $d(a_i,b_j)$ 表示把信源符号 $a_i$ 复现为 $b_j$ 时产生的失真。
\end{definitionBox}

失真矩阵由问题的保真度准则决定，不是由概率分布决定。做题时应先看题目给出的失真矩阵，再考虑试验信道 $p(y\mid x)$ 怎样选择。

\begin{definitionBox}{汉明失真测度}
当输入和输出符号一一对应时，常用汉明失真测度：
\[
d(a_i,b_j)=
\begin{cases}
0, & i=j,\\
1, & i\ne j.
\end{cases}
\]
它只关心复现符号是否正确，不区分不同错误之间的严重程度。
\end{definitionBox}

例如 $n$ 元汉明失真矩阵可写成
\[
\mathbf d=
\begin{bmatrix}
0&1&\cdots&1\\
1&0&\cdots&1\\
\vdots&\vdots&\ddots&\vdots\\
1&1&\cdots&0
\end{bmatrix}.
\]

\subsubsection{序列失真测度}

实际编码通常不是逐个符号独立处理，而是对长度为 $N$ 的序列编码。设
\[
\mathbf x=(x_1,x_2,\dots,x_N),
\qquad
\mathbf y=(y_1,y_2,\dots,y_N),
\]
则序列失真通常定义为各位置单符号失真的平均：
\[
d_N(\mathbf x,\mathbf y)=\frac1N\sum_{i=1}^{N}d(x_i,y_i).
\]

这个定义说明本章关心的是每个信源符号平均承担多少失真。少数位置失真较大并不一定违反要求，关键是整体平均失真是否超过允许值。

\begin{formulaBox}{平均失真}
给定信源分布 $p(x)$ 和试验信道 $p(y\mid x)$，单符号平均失真为
\[
D=E[d(X,Y)]=\sum_{x,y}p(x)p(y\mid x)d(x,y).
\]
对长度为 $N$ 的序列，有
\[
D_N=E[d_N(\mathbf X,\mathbf Y)]
=\frac1N\sum_{i=1}^{N}E[d(X_i,Y_i)].
\]
\end{formulaBox}

平均失真公式是后面所有计算的入口。若题目给出一个确定的映射规则，则 $p(y\mid x)$ 只有 $0$ 和 $1$；若给出随机试验信道，就按条件概率加权求和。

\subsection{信息率失真函数}

\subsubsection{信息率失真函数的定义}

给定允许平均失真 $D$，所有满足
\[
\sum_{x,y}p(x)p(y\mid x)d(x,y)\le D
\]
的试验信道组成集合 $P_D$。在这些试验信道中，$I(X;Y)$ 表示复现端关于信源端至少保留了多少信息。为了在失真不超过 $D$ 的前提下尽量压缩，应让这个互信息尽可能小。

\begin{definitionBox}{信息率失真函数}
离散信源在失真测度 $d(x,y)$ 下的信息率失真函数定义为
\[
R(D)=\min_{p(y\mid x)\in P_D} I(X;Y),
\]
其中 $P_D$ 是满足平均失真约束的试验信道集合。

等价地，
\[
R(D)=\min_{p(y\mid x)\in P_D}
\sum_{x,y}p(x)p(y\mid x)
\log\frac{p(y\mid x)}{\sum_x p(x)p(y\mid x)}.
\]
\end{definitionBox}

若对数以 $2$ 为底，$R(D)$ 的单位为 bit/信源符号。它给出允许平均失真不超过 $D$ 时，每个信源符号至少需要保留的信息量。$D$ 越大，允许压缩越狠，所需码率越低。

\subsubsection{定义域与端点}

信息率失真函数只在可实现的失真范围内讨论。最小可能失真来自对每个输入符号选择失真最小的输出；最大有意义失真来自完全不看输入、始终输出某个固定符号时能达到的最小平均失真。

\begin{formulaBox}{$R(D)$ 的定义域端点}
信息率失真函数的定义域通常写成
\[
D_{\min}\le D\le D_{\max},
\]
其中
\[
D_{\min}=\sum_x p(x)\min_y d(x,y),
\]
\[
D_{\max}=\min_y\sum_x p(x)d(x,y).
\]
端点处满足
\[
R(D_{\min})=H(X)\quad\text{常见无零失真偏移时写作 }R(0)=H(X),
\]
\[
R(D_{\max})=0.
\]
\end{formulaBox}

若失真矩阵对每个输入都有一个零失真复现符号，则 $D_{\min}=0$，这时无失真端点就是 $R(0)=H(X)$。当 $D=D_{\max}$ 时，可以只输出一个固定符号，不再向复现端传递关于 $X$ 的信息，因此 $I(X;Y)=0$，所以 $R(D_{\max})=0$。

\begin{example}
\label{ex:ch9-dmin-dmax}
\textbf{例9.2.1：计算 $D_{\min}$ 与 $D_{\max}$}：设试验信道输入符号为 $\{a_1,a_2,a_3\}$，且
\[
p(a_1)=p(a_2)=p(a_3)=\frac13.
\]
失真矩阵为
\[
(d_{ij})=
\begin{bmatrix}
1&2&3\\
2&1&3\\
3&2&1
\end{bmatrix}.
\]
求 $D_{\min}$、$D_{\max}$ 及对应试验信道转移概率矩阵。

\noindent \textbf{解题流程}：\\[2pt]
步骤一：$D_{\min}$ 对每一行取最小失真，再按 $p(a_i)$ 加权。\\
步骤二：$D_{\max}$ 对每一列求加权平均失真，再取最小列。\\
步骤三：端点对应的试验信道分别是“每行选最小项”和“所有输入都映射到最优固定输出”。

\textbf{解}：由
\[
D_{\min}=\sum_x p(x)\min_y d(x,y)
\]
得
\[
D_{\min}=\frac13\cdot1+\frac13\cdot1+\frac13\cdot1=1.
\]
每一行的最小元素都在对角线上，因此对应转移概率矩阵为
\[
P_{Y\mid X}=
\begin{bmatrix}
1&0&0\\
0&1&0\\
0&0&1
\end{bmatrix}.
\]

再由
\[
D_{\max}=\min_y\sum_x p(x)d(x,y)
\]
分别计算三列平均失真：
\[
\frac13(1+2+3)=2,
\qquad
\frac13(2+1+2)=\frac53,
\qquad
\frac13(3+3+1)=\frac73.
\]
最小值为第二列，所以
\[
D_{\max}=\frac53.
\]
对应做法是不管输入是什么，都输出 $b_2$：
\[
P_{Y\mid X}=
\begin{bmatrix}
0&1&0\\
0&1&0\\
0&1&0
\end{bmatrix}.
\]
\end{example}

\subsubsection{信息率失真函数的性质}

信息率失真函数的形状可以帮助判断计算结果是否合理。随着允许失真增大，所需码率不应上升；同时，因为可以在不同试验信道之间随机切换，$R(D)$ 具有凸性。

\begin{theoremBox}{$R(D)$ 的基本性质}
在定义域内，信息率失真函数满足：
\begin{enumerate}
  \item $R(D)$ 是关于 $D$ 的下凸函数，即对 $0<\alpha<1$，
  \[
  R\bigl(\alpha D_1+(1-\alpha)D_2\bigr)
  \le \alpha R(D_1)+(1-\alpha)R(D_2).
  \]
  \item $R(D)$ 在 $(D_{\min},D_{\max})$ 上连续。
  \item $R(D)$ 在有效区间内单调非增，通常表现为严格递减。
\end{enumerate}
\end{theoremBox}

单调性的直觉很直接：若 $D_1>D_2$，则允许失真为 $D_1$ 的试验信道集合包含允许失真为 $D_2$ 的集合，即 $P_{D_1}\supset P_{D_2}$。在更大的集合中取最小互信息，不可能比在更小集合中取到的值更大，所以 $R(D_1)\le R(D_2)$。

下凸性对应“时间共享”思想：一部分序列使用失真 $D_1$、码率 $R(D_1)$ 的方案，另一部分使用失真 $D_2$、码率 $R(D_2)$ 的方案，整体平均失真和码率就是二者的凸组合。最优 $R(D)$ 不会比这种混合方案更差。

\subsubsection{二元等概信源的汉明失真}

PPT 中用二元信源曲线说明 $R(D)$ 随信源分布变化的形状。考试中最常见的是二元等概信源配汉明失真，此时公式特别简单。

\begin{formulaBox}{二元等概信源的率失真函数}
若二元信源 $X\in\{0,1\}$ 等概分布，失真测度为汉明失真，则
\[
R(D)=1-H(D),
\qquad 0\le D\le\frac12,
\]
其中
\[
H(D)=-D\log_2D-(1-D)\log_2(1-D).
\]
当 $D\ge1/2$ 时，$R(D)=0$。
\end{formulaBox}

这个公式说明：若完全不允许错误，$D=0$，则 $R(0)=1$ bit/符号；若允许平均错误率达到 $1/2$，直接固定输出 $0$ 或 $1$ 已经足够，因而不需要传输任何信息。

\subsection{限失真信源编码定理}

\subsubsection{码率压缩的含义}

设信源发出长度为 $N$ 的序列，编码器只准备 $M$ 个码字。若要求无失真编码，典型序列数大约为 $2^{NH}$，因此需要
\[
M>2^{NH}
\]
才可能在 $N$ 足够大时实现无失真恢复。令
\[
R=\frac{\log_2 M}{N},
\]
则无失真编码要求 $R>H$。

限失真编码的思想是：当 $R<H$ 时，不再试图区分所有可能信源序列，而是把多个相近序列归并到同一个复现码字。编码器按最小失真原则为输入序列寻找一个码字 $\mathbf y_m$，只发送码字序号 $m$；译码器收到序号后输出对应复现序列。这样用更少的码字覆盖大量源序列，代价是产生可控失真。

\begin{example}
\label{ex:ch9-single-symbol-compression}
\textbf{例9.3.1：单符号压缩算法}：设信源 $X$ 的符号集为
\[
\{a_1,a_2,\dots,a_{2n}\},
\]
且等概分布 $p_i=1/(2n)$。失真测度为汉明测度
\[
d_{ij}=
\begin{cases}
1, & i\ne j,\\
0, & i=j.
\end{cases}
\]
设计一种单符号压缩算法，使平均失真 $D=1/2$，并求压缩后的码率 $R$。

\noindent \textbf{解题流程}：\\[2pt]
步骤一：把 $2n$ 个输入符号压缩到 $n$ 个输出符号。\\
步骤二：按映射规则统计哪些输入会产生汉明失真 $1$。\\
步骤三：求输出分布熵，作为该单符号压缩方案的平均码率。

\textbf{解}：设压缩输出符号集为
\[
\{b_1,b_2,\dots,b_n\}.
\]
采用映射
\[
a_i\to b_i,
\qquad i=1,2,\dots,n-1,
\]
并令
\[
a_j\to b_n,
\qquad j=n,n+1,\dots,2n.
\]
只有 $a_{n+1},\dots,a_{2n}$ 被映射到与自身不同的复现符号，因此
\[
D=\sum_{x,y}p(x)p(y\mid x)d(x,y)
=\frac{1}{2n}\sum_{j=n+1}^{2n}1
=\frac12.
\]
输出分布为
\[
p(b_j)=\frac{1}{2n},\quad j=1,2,\dots,n-1,
\]
\[
p(b_n)=\frac{n+1}{2n}.
\]
所以压缩后码率为输出熵
\[
R=H(Y)=\frac{n-1}{2n}\log_2(2n)
+\frac{n+1}{2n}\log_2\frac{2n}{n+1}.
\]
整理得
\[
R=\log_2(2n)-\frac{n+1}{2n}\log_2(n+1).
\]
\end{example}

\subsubsection{香农第三定理}

率失真函数不仅是一个优化定义，还是真正的编码理论下界。香农第三定理说明，只要码长足够长，码率略高于 $R(D)$ 就能达到平均失真 $D$；若码率低于 $R(D)$，则不可能达到该失真要求。

\begin{theoremBox}{定理9.3.1：限失真信源编码定理（香农第三定理）}
任意给定 $\varepsilon>0$，总存在一种信源编码，使当
\[
R\ge R(D)+\varepsilon
\]
时，平均失真满足
\[
\bar D\le D+\varepsilon.
\]
反之，如果
\[
R<R(D),
\]
则不可能存在使平均失真
\[
\bar D\le D
\]
的编码。
\end{theoremBox}

这个定理和前两条香农定理形式一致：香农第一定理给出无失真压缩下界 $H$，香农第二定理给出有噪信道可靠传输上界 $C$，香农第三定理给出允许失真 $D$ 时的压缩下界 $R(D)$。它同样是存在性定理，不直接给出具体最优编码构造。

\subsection{限失真信源信道编码定理}

如果压缩后的信息还要经过有噪信道传输，就需要把率失真函数和信道容量联系起来。直觉上，信道每秒能可靠承载的信息量必须不小于信源在失真要求 $D$ 下每秒至少需要的信息量。

\begin{theoremBox}{定理9.3.2：限失真信源信道编码定理}
设离散无记忆信源的信息率失真函数为 $R(D)$，离散无记忆信道容量为 $C$。若
\[
C\ge R(D),
\]
则存在编码系统，使信源序列通过信道传输后的平均失真不超过 $D$。

若
\[
C<R(D),
\]
则不能使传输后的平均失真不超过 $D$。
\end{theoremBox}

使用该定理时要注意单位一致。若信源每秒发出 $r_s$ 个符号，而 $R(D)$ 的单位是 bit/信源符号，则每秒所需传输率为 $r_sR(D)$。它应与信道每秒容量 $C$ 比较。

\begin{example}
\label{ex:ch9-binary-source-distortion-transmission}
\textbf{习题9.25：二元信源的限失真传输}：有一信源
\[
\begin{pmatrix}
S\\ P
\end{pmatrix}
=
\begin{pmatrix}
s_1&s_2\\
1/2&1/2
\end{pmatrix},
\]
每秒发出 $2.66$ 个符号。将信源输出通过二元无噪信道传输，而信道每秒仅传送两个二元符号。问能否无失真传输？若不能，在汉明失真下允许多大平均失真便可以通过该信道传输？

\noindent \textbf{解题流程}：\\[2pt]
步骤一：先比较信源每秒熵和信道每秒容量，判断无失真传输是否可能。\\
步骤二：若不能无失真传输，用二元等概汉明失真公式 $R(D)=1-H(D)$。\\
步骤三：令信源每秒所需率不超过信道容量，解出最小允许失真。

\textbf{解}：信源每个符号熵为
\[
H(X)=1\text{ bit/符号}.
\]
每秒发出 $2.66$ 个符号，所以无失真传输所需信息率为
\[
2.66\text{ bit/s}.
\]
二元无噪信道每个信道符号容量为 $1$ bit，每秒传送 $2$ 个二元符号，所以
\[
C=2\text{ bit/s}.
\]
由于
\[
2.66>C,
\]
所以不能无失真传输。

在汉明失真下，二元等概信源满足
\[
R(D)=1-H(D).
\]
限失真传输要求
\[
2.66R(D)\le2,
\]
即
\[
R(D)\le\frac{2}{2.66}\approx0.7519.
\]
因此
\[
1-H(D)\le0.7519,
\qquad
H(D)\ge0.2481.
\]
查二元熵函数或数值求解得
\[
D\ge0.0413.
\]
所以平均失真至少允许约 $0.0413$，才可能通过该信道完成限失真传输。
\end{example}

\subsection{题型整理：第九章常见计算}

\textbf{题型快速判断}：看到题目后先判断它是在算失真端点、平均失真、率失真函数，还是在比较信源所需率和信道容量。
\begin{itemize}
  \item 给失真矩阵和信源分布 $\to$ 常先算 $D_{\min}$、$D_{\max}$；
  \item 给具体映射或试验信道 $p(y\mid x)$ $\to$ 用 $D=\sum p(x)p(y\mid x)d(x,y)$ 算平均失真；
  \item 问 $R(D)$ 的定义或性质 $\to$ 写 $R(D)=\min_{p(y\mid x)\in P_D}I(X;Y)$，并说明单调非增、下凸、连续；
  \item 二元等概信源加汉明失真 $\to$ 优先用 $R(D)=1-H(D)$，并记得适用区间 $0\le D\le1/2$；
  \item 给信源符号率和信道容量 $\to$ 比较 $r_sR(D)$ 与 $C$，无失真时比较 $r_sH(X)$ 与 $C$。
\end{itemize}

\textbf{答题模板}：
\begin{enumerate}
  \item \textbf{端点计算题}：$D_{\min}$ 按行取最小后加权；$D_{\max}$ 按列求加权平均后取最小。
  \item \textbf{平均失真题}：先写清 $p(y\mid x)$，再把非零映射对应的失真逐项加权。
  \item \textbf{定理判断题}：若是限失真信源编码，比较 $R$ 与 $R(D)$；若通过信道传输，比较 $C$ 与信源每秒所需的率。
  \item \textbf{二元汉明题}：把 $R(D)=1-H(D)$ 代入容量约束，解出 $D$；若只问端点，则 $D_{\min}=0$、$D_{\max}=1/2$、$R(D_{\max})=0$。
\end{enumerate}

\begin{warningBox}{第九章易错点汇总}
\begin{itemize}
  \item $R(D)$ 是在满足失真约束的试验信道中\textbf{最小化} $I(X;Y)$，不要和信道容量的最大化混淆。
  \item $D_{\min}$ 是每一行先取最小再按输入概率加权；$D_{\max}$ 是每一列先按输入概率加权再取最小，二者顺序不同。
  \item $R(D_{\max})=0$ 的含义是可以固定输出一个再现符号，不是说失真为零。
  \item 二元等概汉明失真公式 $R(D)=1-H(D)$ 只在 $0\le D\le1/2$ 内使用；当 $D\ge1/2$ 时 $R(D)=0$。
  \item 比较信源和信道时必须统一单位：bit/符号要乘以符号率后才能和 bit/s 的信道容量比较。
  \item 香农第三定理是存在性结论，不表示任意有限码长或任意具体编码都能恰好达到 $R(D)$。
\end{itemize}
\end{warningBox}

\subsection{公式速查表}

{\scriptsize
\noindent\begin{tabularx}{\textwidth}{@{}l>{\raggedright\arraybackslash}X@{}}
\toprule
\textbf{类别} & \textbf{公式 / 核心结论} \\
\midrule
失真矩阵 & $\mathbf d=[d(a_i,b_j)]$，元素表示 $a_i$ 复现为 $b_j$ 的失真 \\
\midrule
汉明失真 & \textcolor{Red}{$d(a_i,b_j)=0$ 若 $i=j$，$d(a_i,b_j)=1$ 若 $i\ne j$} \\
\midrule
序列失真 & \textcolor{Red}{$d_N(\mathbf x,\mathbf y)=\dfrac1N\sum_{i=1}^{N}d(x_i,y_i)$} \\
\midrule
平均失真 & \textcolor{Red}{$D=\sum_{x,y}p(x)p(y\mid x)d(x,y)$} \\
\midrule
率失真函数 & \textcolor{Red}{$R(D)=\min_{p(y\mid x)\in P_D}I(X;Y)$}，其中 $P_D$ 满足平均失真不超过 $D$ \\
\midrule
互信息展开 & $R(D)=\min\sum_{x,y}p(x)p(y\mid x)\log\dfrac{p(y\mid x)}{\sum_xp(x)p(y\mid x)}$ \\
\midrule
最小失真端点 & \textcolor{Red}{$D_{\min}=\sum_xp(x)\min_y d(x,y)$} \\
\midrule
最大有意义失真 & \textcolor{Red}{$D_{\max}=\min_y\sum_xp(x)d(x,y)$}，且 $R(D_{\max})=0$ \\
\midrule
下凸性 & $R(\alpha D_1+(1-\alpha)D_2)\le\alpha R(D_1)+(1-\alpha)R(D_2)$ \\
\midrule
二元等概汉明 & \textcolor{Red}{$R(D)=1-H(D)$，$0\le D\le1/2$；$D\ge1/2$ 时 $R(D)=0$} \\
\midrule
限失真信源编码 & \textcolor{Red}{$R\ge R(D)+\varepsilon$ 时可使 $\bar D\le D+\varepsilon$；$R<R(D)$ 时不能使 $\bar D\le D$} \\
\midrule
限失真信源信道编码 & \textcolor{Red}{$C\ge R(D)$ 时平均失真可不超过 $D$；$C<R(D)$ 时不能达到该失真要求} \\
\bottomrule
\end{tabularx}
}

\subsection{本章要求}

本章学习重点：
\begin{itemize}
  \item 理解有失真信源编码模型和试验信道 $p(y\mid x)$ 的含义；
  \item 掌握单符号失真测度、失真矩阵和汉明失真测度；
  \item 会写序列失真 $d_N(\mathbf x,\mathbf y)$ 和平均失真 $D$；
  \item 掌握信息率失真函数 $R(D)=\min I(X;Y)$ 的定义和物理意义；
  \item 会计算 $D_{\min}$、$D_{\max}$ 及对应的试验信道转移概率矩阵；
  \item 理解 $R(D)$ 的下凸性、连续性和单调递减性质；
  \item 掌握二元等概信源在汉明失真下的公式 $R(D)=1-H(D)$；
  \item 能完成例9.3.1 类型的简单压缩映射、平均失真和码率计算；
  \item 理解限失真信源编码定理，即香农第三定理；
  \item 理解限失真信源信道编码定理，并会统一 bit/符号与 bit/s 后比较 $R(D)$ 和 $C$；
  \item 会用习题9.25 类型的方法判断是否能无失真传输，以及求最小允许平均失真。
\end{itemize}
